\subsection{Haces sobre un espacio localmente anillado}
\label{subsection:0.5.5-fr}

\begin{env}[5.5.1]
\label{0.5.5.1-fr}
Diremos que un espacio anillado $(X,\mathcal{O}_X)$ es un \emph{espacio
anillado en anillos locales} si, para todo $x\in X$, $\mathcal{O}_x$ es un
anillo local; estos espacios anillados serán con mucho los más frecuentes de los
que se considerarán en este trabajo. Designaremos entonces por
$\mathfrak{m}_x$ el \emph{ideal maximal} de $\mathcal{O}_x$, por $\boldsymbol{k}(x)$ el
\emph{cuerpo residual} $\mathcal{O}_x/\mathfrak{m}_x$; para todo
$\mathcal{O}_X$-Módulo $\mathcal{F}$, todo abierto $U$ de $X$, todo punto
$x\in U$ y toda sección $f\in\Gamma(U,\mathcal{F})$, designaremos por
$f(x)$ la \emph{clase} del germen
$f_x\in\mathcal{F}_x$ mod.~$\mathfrak{m}_x\mathcal{F}_x$, y diremos que
es el \emph{valor} de $f$ en el punto $x$. La relación $f(x)=0$ significa,
pues, que $f_x\in\mathfrak{m}_x\mathcal{F}_x$; cuando se cumple, se dirá
(por abuso de lenguaje) que $f$ \emph{se anula en $x$}. Se tendrá cuidado de no
confundir esta relación con $f_x=0$.
\end{env}

\begin{env}[5.5.2]
\label{0.5.5.2-fr}
Sean $X$ un espacio localmente anillado, $\mathcal{L}$ un
$\mathcal{O}_X$-Módulo invertible, $f$ una sección de $\mathcal{L}$ sobre
$X$. Hay entonces \emph{equivalencia} entre las tres propiedades siguientes
en un punto $x\in X$:
\begin{enumerate}
  \item[a)] \emph{$f_x$ es un generador de $\mathcal{L}_x$};
  \item[b)] \emph{$f_x\notin\mathfrak{m}_x\mathcal{L}_x$} (es decir,
    $f(x)\neq0$);
  \item[c)] \emph{Existe una sección $g$ de $\mathcal{L}^{-1}$ sobre
    una vecindad abierta $V$ de $x$ tal que la imagen canónica de $f\otimes
    g$ en $\Gamma(V,\mathcal{O}_X)$ \textup{(5.4.3)} es la sección unidad}.
\end{enumerate}

En efecto, siendo la cuestión local, se puede limitar al caso en que
$\mathcal{L}=\mathcal{O}_X$; la equivalencia de $a)$ y $b)$ es entonces
evidente, y es claro que $c)$ implica $b)$. Por otra parte, si
$f_x\notin\mathfrak{m}_x$, $f_x$ es invertible en $\mathcal{O}_x$, sea
$f_xg_x=1_x$. Por definición de los gérmenes de secciones, esto quiere decir que
existe una vecindad $V$ de $x$ y una sección $g$ de $\mathcal{O}_X$ sobre
$V$ tal que $fg=1$ en $V$, de donde $c)$.

Resulta inmediatamente de la condición $c)$ que el conjunto $X_f$ de los $x$
que verifican las condiciones equivalentes $a)$, $b)$, $c)$ es \emph{abierto} en
$X$; según la terminología introducida en (5.5.1), es el conjunto de los
$x$ donde $f$ \emph{no se anula}.
\end{env}

\begin{env}[5.5.3]
\label{0.5.5.3-fr}
Bajo las hipótesis de (5.5.2), sea $\mathcal{L}'$ un segundo
$\mathcal{O}_X$-Módulo invertible; entonces, si
$f\in\Gamma(X,\mathcal{L})$, $g\in\Gamma(X,\mathcal{L}')$, se tiene
\[
  X_f\cap X_g=X_{f\otimes g}
\]

En efecto, se reduce inmediatamente al caso en que
$\mathcal{L}=\mathcal{L}'=\mathcal{O}_X$ (siendo la cuestión local); como
$f\otimes g$ se identifica entonces canónicamente con el producto $fg$, la proposición
es evidente.
\end{env}

\oldpage[0\textsubscript{I}]{54}

\begin{env}[5.5.4]
\label{0.5.5.4-fr}
Sea $\mathcal{F}$ un $\mathcal{O}_X$-Módulo localmente libre de rango $n$;
es inmediato que $\bigwedge^p\mathcal{F}$ es un $\mathcal{O}_X$-Módulo
localmente libre de rango $\binom{n}{p}$ si $p\leq n$, reducido a $0$ si $p>n$,
pues la cuestión es local y se reduce al caso en que
$\mathcal{F}=\mathcal{O}_X^n$; además, para todo $x\in X$,
\[
  (\bigwedge^p\mathcal{F})_x/
  \mathfrak{m}_x(\bigwedge^p\mathcal{F})_x
\]
es un espacio vectorial de dimensión $\binom{n}{p}$ sobre $\boldsymbol{k}(x)$, que
se identifica canónicamente con
\[
  \bigwedge^p(\mathcal{F}_x/\mathfrak{m}_x\mathcal{F}_x).
\]
Sean
$s_1,\ldots,s_p$ unas secciones de $F$ sobre un abierto $U$ de $X$, y
sea $s=s_1\wedge\ldots\wedge s_p$, que es una sección de
$\bigwedge^p\mathcal{F}$ sobre $U$ (4.1.5); se tiene
$s(x)=s_1(x)\wedge\ldots\wedge s_p(x)$, y, por consiguiente, decir que
$s_1(x),\ldots,s_p(x)$ son \emph{linealmente dependientes} significa que
$s(x)=0$. Se concluye que el \emph{conjunto de los $x\in X$ tales que
$s_1(x),\ldots,s_p(x)$ son linealmente independientes es abierto en
$X$}: basta, en efecto, reduciéndose al caso en que
$\mathcal{F}=\mathcal{O}_X^n$, aplicar (5.5.2) a la sección imagen de $s$
por una de las proyecciones de
$\bigwedge^p\mathcal{F}=\mathcal{O}_X^{\binom{n}{p}}$ sobre los
$\binom{n}{p}$ factores.

En particular, si $s_1,\ldots,s_n$ son $n$ secciones de $\mathcal{F}$
sobre $U$ tales que $s_1(x),\ldots,s_n(x)$ son linealmente
independientes en todo punto $x\in U$, el homomorfismo
$u:\mathcal{O}_X^n|U\to\mathcal{F}|U$ definido por los $s_i$ (5.1.1) es un
\emph{isomorfismo}: en efecto, se puede restringir al caso en que
$\mathcal{F}=\mathcal{O}_X^n$ y en que se identifican canónicamente
$\bigwedge^n\mathcal{F}$ y $\mathcal{O}_X$;
$s=s_1\wedge\ldots\wedge s_n$ es entonces una sección de $\mathcal{O}_X$
\emph{invertible} sobre $U$, y se define un homomorfismo inverso
de $u$ mediante las fórmulas de Cramer.
\end{env}

\begin{env}[5.5.5]
\label{0.5.5.5-fr}
Sean $\mathcal{E}$, $\mathcal{F}$ dos $\mathcal{O}_X$-Módulos localmente
libres (de rango finito), y sea $u:\mathcal{E}\to\mathcal{F}$ un
homomorfismo. Para que exista una vecindad $u$ de $x\in X$ tal que $u|U$
sea \emph{inyectivo} y $\mathcal{F}|U$ sea \emph{suma directa de
$u(\mathcal{E})|U$ y de un sub-$(\mathcal{O}_X|U)$-Módulo localmente libre
$\mathcal{G}$}, es necesario y suficiente que
$u_x:\mathcal{E}_x\to\mathcal{F}_x$ dé, por paso a los cocientes, un
homomorfismo \emph{inyectivo} de espacios vectoriales
\[
  \mathcal{E}_x/\mathfrak{m}_x\mathcal{E}_x
  \longrightarrow
  \mathcal{F}_x/\mathfrak{m}_x\mathcal{F}_x.
\]
La condición es, en efecto, \emph{necesaria}, pues $\mathcal{F}_x$ es entonces
suma directa de los $\mathcal{O}_x$-módulos libres $u_x(\mathcal{E}_x)$ y
$\mathcal{G}_x$, por tanto $\mathcal{F}_x/\mathfrak{m}_x\mathcal{F}_x$ es suma
directa de
$u_x(\mathcal{E}_x)/\mathfrak{m}_x u_x(\mathcal{E}_x)$ y de
$\mathcal{G}_x/\mathfrak{m}_x\mathcal{G}_x$. La condición es
\emph{suficiente}, pues se puede limitar al caso en que
$\mathcal{E}=\mathcal{O}_X^m$; sean $s_1,\ldots,s_m$ las imágenes por $u$ de las
secciones $e_i$ de $\mathcal{O}_X^m$ tales que $(e_i)_y$ es igual al $i$-ésimo
elemento de la base canónica de $\mathcal{O}_y^m$ para todo $y\in X$
(\emph{secciones canónicas} de $\mathcal{O}_X^m$); por hipótesis,
$s_1(x),\ldots,s_m(x)$ son linealmente independientes, por lo que, si $\mathcal{F}$
es de rango $n$, existen $n-m$ secciones $s_{m+1},\ldots,s_n$ de
$\mathcal{F}$ sobre una vecindad $V$ de $x$ tales que los $s_i(x)$
$(1\leq i\leq n)$ forman una base de
$\mathcal{F}_x/\mathfrak{m}_x\mathcal{F}_x$. Existe entonces (5.5.4) una
vecindad $U\subset V$ de $x$ tal que los $s_i(y)$ $(1\leq i\leq n)$ forman
una base de $\mathcal{F}_y/\mathfrak{m}_y\mathcal{F}_y$ para todo $y\in V$,
y se concluye (5.5.4) que hay un isomorfismo de $\mathcal{F}|U$ sobre
$\mathcal{O}_X^n|U$, que aplica los $s_i|U$ $(1\leq i\leq m)$ en los
$e_i|U$, lo que termina la demostración.
\end{env}

